\ifdefined\gkver\else\def\gkver{student}\fi
\documentclass[\gkver]{gaokaozhenti}
\begin{document}

\chapter{2024年广东}

\begin{enumerate}
\item
%% number: 1
%% paperName: 2024年广东高考第 1 题 4 分
%% typeId: 1
%% type: 单选题
%% chapter: 交流电
%% point: 交流电
%% method: 无
%% score: 4
%% degree: 850
%% duplicateId: 0
%% body:
将阻值为 $50 \UO$ 的电阻接在正弦式交流电源上。电阻两端电压随时间的变化规律如图所示。下列说法正确的是\xzanswer{D}

\onepicture[w=5.00cm, align=right]{figs/01.png}

\fourchoices[answer=D]
{该交流电的频率为 $100 \UHz$}
{通过电阻电流的峰值为 $0.2 \UA$}
{电阻在 1 秒内消耗的电能为 $1 \UJ$}
{电阻两端电压表达式为 $u = 10\sqrt{2}\sin(100\pi t) \UV$}

%% answer: D
%% memo:
\memoanswer{
A．由图可知交流电的周期为 $0.02 \Us$，则频率为 $f = \frac{1}{T} = 50 \UHz$。故 A 错误；

B．根据图像可知电压的峰值为 $10\sqrt{2} \UV$，根据欧姆定律可知电流的峰值

$I_{m}=\frac{U_{m}}{R}=\frac{10\sqrt{2} \UV}{50 \UO}=0.2\sqrt{2} \UA$

故 B 错误；

C．电流的有效值为 $I = \frac{I_{m}}{\sqrt{2}} = 0.2 \UA$

所以电阻在 $1 \Us$ 内消耗的电能为 $W = I^{2}Rt = 0.2^{2} \times 50 \times 1 \UJ = 2 \UJ$

故 C 错误；

D．根据图像可知其电压表达式为

$u = U_{m}\sin\omega t = 10\sqrt{2}\sin\frac{2\pi}{T}t \left( \UV \right) = 10\sqrt{2}\sin100\pi t \left( \UV \right)$

故 D 正确。

故选 D。
}

\item
%% number: 2
%% paperName: 2024年广东高考第 2 题 4 分
%% typeId: 1
%% type: 单选题
%% chapter: 原子和原子核
%% point: 人工核反应
%% method: 无
%% score: 4
%% degree: 850
%% duplicateId: 0
%% body:
我国正在建设的大科学装置——“强流重离子加速器”。其科学目标之一是探寻神秘的“119 号”元素，科学家尝试使用核反应 \ce{Y + ^{243}_{95}Am -> ^{A}_{119}X + 2^{1}_{0}n} 产生该元素。关于原子核 Y 和质量数 $A$，下列选项正确的是\xzanswer{C}

\fourchoices[answer=C]
{Y 为 \ce{^{58}_{26}Fe}，$A = 299$}
{Y 为 \ce{^{58}_{26}Fe}，$A = 301$}
{Y 为 \ce{^{54}_{24}Cr}，$A = 295$}
{Y 为 \ce{^{54}_{24}Cr}，$A = 297$}

%% answer: C
%% memo:
\memoanswer{
根据电荷数守恒设 Y 的质子数为 $y$，则有

$y + 95 = 119 + 0$

可得 $y = 24$

即 Y 为 \ce{^{54}_{24}Cr}；根据质量数守恒，则有

$54 + 243 = A + 2$

可得 $A = 295$

故选 C。
}

\item
%% number: 3
%% paperName: 2024年广东高考第 3 题 4 分
%% typeId: 1
%% type: 单选题
%% chapter: 机械波
%% point: 波的图象
%% method: 无
%% score: 4
%% degree: 850
%% duplicateId: 0
%% body:
一列简谐横波沿 $x$ 轴正方向传播，波速为 $1 \Ums$，$t = 0$ 时的波形如图所示。$t = 1 \Us$ 时，$x = 1.5 \Um$ 处的质点相对平衡位置的位移为\xzanswer{B}

\onepicture[w=5.00cm, align=right]{figs/03.png}

\fourchoices[answer=B]
{$0$}
{$0.1 \Um$}
{$- 0.1 \Um$}
{$0.2 \Um$}

%% answer: B
%% memo:
\memoanswer{
由图可知简谐波的波长为 $\lambda = 2 \Um$，所以周期为 $T = \frac{\lambda}{v} = 2 \Us$。

当 $t = 1 \Us$ 时，$x = 1.5 \Um$ 处的质点运动半个周期到达波峰处，故相对平衡位置的位移为 $0.1 \Um$。

故选 B。
}

\item
%% number: 4
%% paperName: 2024年广东高考第 4 题 4 分
%% typeId: 1
%% type: 单选题
%% chapter: 电磁感应
%% point: 楞次定律
%% method: 无
%% score: 4
%% degree: 650
%% duplicateId: 0
%% body:
电磁俘能器可在汽车发动机振动时利用电磁感应发电实现能量回收，结构如图甲所示。两对永磁铁可随发动机一起上下振动，每对永磁铁间有水平方向的匀强磁场，磁感应强度大小均为 $B$。磁场中，边长为 $L$ 的正方形线圈竖直固定在减震装置上。某时刻磁场分布与线圈位置如图乙所示，永磁铁振动时磁场分界线不会离开线圈。关于图乙中的线圈，下列说法正确的是\xzanswer{D}

\onepicture[w=6.00cm]{figs/04.png}

\fourchoices[answer=D]
{穿过线圈的磁通量为 $BL^{2}$}
{永磁铁相对线圈上升越高，线圈中感应电动势越大}
{永磁铁相对线圈上升越快，线圈中感应电动势越小}
{永磁铁相对线圈下降时，线圈中感应电流的方向为顺时针方向}

%% answer: D
%% memo:
\memoanswer{
A．根据图乙可知此时穿过线圈的磁通量为 0，故 A 错误；

BC．根据法拉第电磁感应定律可知永磁铁相对线圈上升越快，磁通量变化越快，线圈中感应电动势越大，故 BC 错误；

D．永磁铁相对线圈下降时，根据楞次定律可知线圈中感应电流的方向为顺时针方向，故 D 正确。

故选 D。
}

\item
%% number: 5
%% paperName: 2024年广东高考第 5 题 4 分
%% typeId: 1
%% type: 单选题
%% chapter: 曲线运动
%% point: 圆周运动的应用
%% method: 无
%% score: 4
%% degree: 550
%% duplicateId: 0
%% body:
如图所示，在细绳的拉动下，半径为 $r$ 的卷轴可绕其固定的中心点 O 在水平面内转动。卷轴上沿半径方向固定着长度为 $l$ 的细管，管底在 O 点。细管内有一根原长为 $\frac{l}{2}$、劲度系数为 $k$ 的轻质弹簧，弹簧底端固定在管底，顶端连接质量为 $m$、可视为质点的插销。当以速度 $v$ 匀速拉动细绳时，插销做匀速圆周运动。若 $v$ 过大，插销会卡进固定的端盖，使卷轴转动停止。忽略摩擦力，弹簧在弹性限度内。要使卷轴转动不停止，$v$ 的最大值为\xzanswer{A}

\onepicture[w=4.50cm, align=right]{figs/05.png}

\fourchoices[answer=A]
{$r\sqrt{\frac{k}{2m}}$}
{$l\sqrt{\frac{k}{2m}}$}
{$r\sqrt{\frac{2k}{m}}$}
{$l\sqrt{\frac{2k}{m}}$}

%% answer: A
%% memo:
\memoanswer{
由题意可知当插销刚卡紧固定端盖时弹簧的伸长量为 $\Delta x = \frac{l}{2}$，根据胡克定律有

$F = k\Delta x = \frac{kl}{2}$

插销与卷轴同轴转动，角速度相同，对插销有弹力提供向心力 $F = ml\omega^{2}$

对卷轴有 $v = r\omega$

联立解得 $v = r\sqrt{\frac{k}{2m}}$

故选 A。
}

\item
%% number: 6
%% paperName: 2024年广东高考第 6 题 4 分
%% typeId: 1
%% type: 单选题
%% chapter: 光学
%% point: 全反射
%% method: 无
%% score: 4
%% degree: 650
%% duplicateId: 0
%% body:
如图所示，红绿两束单色光，同时从空气中沿同一路径以 $\theta$ 角从 MN 面射入某长方体透明均匀介质。折射光束在 NP 面发生全反射，反射光射向 PQ 面。若 $\theta$ 逐渐增大，两束光在 NP 面上的全反射现象会先后消失。已知在该介质中红光的折射率小于绿光的折射率。下列说法正确的是\xzanswer{B}

\onepicture[w=5.00cm, align=right]{figs/06.png}

\fourchoices[answer=B]
{在 PQ 面上，红光比绿光更靠近 P 点}
{$\theta$ 逐渐增大时，红光的全反射现象先消失}
{$\theta$ 逐渐增大时，入射光可能在 MN 面发生全反射}
{$\theta$ 逐渐减小时，两束光在 MN 面折射的折射角逐渐增大}

%% answer: B
%% memo:
\memoanswer{
A．红光的频率比绿光的频率小，则红光的折射率小于绿光的折射率，在 MN 面，入射角相同，根据折射定律 $n=\frac{\sin\theta}{\sin\alpha}$ 可知绿光在 MN 面的折射角较小，根据几何关系可知绿光比红光更靠近 P 点，故 A 错误；

B．根据全反射发生的条件 $\sin C = \frac{1}{n}$ 可知红光发生全反射的临界角较大，$\theta$ 逐渐增大时，折射光线与 NP 面的交点左移过程中，在 NP 面的入射角先大于红光发生全反射的临界角，所以红光的全反射现象先消失，故 B 正确；

C．在 MN 面，光是从光疏介质到光密介质，无论 $\theta$ 多大，在 MN 面都不可能发生全反射，故 C 错误；

D．根据折射定律 $n = \frac{\sin\theta}{\sin\alpha}$ 可知 $\theta$ 逐渐减小时，两束光在 MN 面折射的折射角逐渐减小，故 D 错误。

故选 B。
}

\item
%% number: 7
%% paperName: 2024年广东高考第 7 题 4 分
%% typeId: 1
%% type: 单选题
%% chapter: 运动和力
%% point: 变加速直线运动
%% method: 无
%% score: 4
%% degree: 450
%% duplicateId: 0
%% body:
如图所示，轻质弹簧竖直放置，下端固定。木块从弹簧正上方 $H$ 高度处由静止释放。以木块释放点为原点，取竖直向下为正方向。木块的位移为 $y$，所受合外力为 $F$，运动时间为 $t$。忽略空气阻力，弹簧在弹性限度内。关于木块从释放到第一次回到原点的过程中，其 $F - y$ 图像或 $y - t$ 图像可能正确的是\xzanswer{B}

\onepicture[w=3.00cm, align=right]{figs/07a.png}

\fourchoices[answer=B, ispicture=true, h=2.3cm]
{figs/07b1.png}
{figs/07b2.png}
{figs/07b3.png}
{figs/07b4.png}

%% answer: B
%% memo:
\memoanswer{
AB．在木块下落 $H$ 高度之前，木块所受合外力为木块的重力保持不变，即

$F = mg$

当木块接触弹簧后到合力为零前，根据牛顿第二定律

$mg - ky = F$

随着 $y$ 增大 $F$ 减小；

当弹簧弹力大于木块的重力后到最低点过程中

$F = ky - mg$

木块所受合外力向上，随着 $y$ 增大 $F$ 增大；$F - y$ 图像如题图 B 所示。

故 B 正确，A 错误；

CD．同理，在木块下落 $H$ 高度之前，木块做匀加速直线运动，根据

$y = \frac{1}{2}gt^{2}$

速度逐渐增大，所以 $y - t$ 图像斜率逐渐增大，当木块接触弹簧后到合力为零前，根据牛顿第二定律

$mg - ky = F$

木块的速度继续增大，做加速度减小的加速运动，所以 $y - t$ 图像斜率继续增大，当弹簧弹力大于木块的重力后到最低点过程中

$F = ky - mg$

木块所受合外力向上，木块做加速度增大的减速运动，所以 $y - t$ 图斜率减小，到达最低点后，木块向上运动，经以上分析可知，木块先做加速度减小的加速运动，再做加速度增大的减速运动，再做匀减速直线运动到最高点，$y - t$ 图像大致为右图。

\onepicture[w=4.00cm]{figs/07_an.jpg}

故 CD 错误。

故选 B。
}

\item
%% number: 8
%% paperName: 2024年广东高考第 8 题 6 分
%% typeId: 2
%% type: 多选题
%% chapter: 电场
%% point: 电场线
%% method: 无
%% score: 6
%% degree: 650
%% duplicateId: 0
%% body:
污水中的污泥絮体经处理后带负电，可利用电泳技术对其进行沉淀去污，基本原理如图所示。涂有绝缘层的金属圆盘和金属棒分别接电源正、负极，金属圆盘置于底部、金属棒插入污水中，形成如图所示的电场分布，其中实线为电场线，虚线为等势面。M 点和 N 点在同一电场线上，M 点和 P 点在同一等势面上。下列说法正确的有\xzanswer{AC}

\onepicture[w=5.00cm, align=right]{figs/08.png}

\fourchoices[answer=AC]
{M 点的电势比 N 点的低}
{N 点的电场强度比 P 点的大}
{污泥絮体从 M 点移到 N 点，电场力对其做正功}
{污泥絮体在 N 点的电势能比其在 P 点的大}

%% answer: AC
%% memo:
\memoanswer{
AC．根据沿着电场线方向电势降低可知 M 点的电势比 N 点的低，污泥絮体带负电，根据 $E_{p} = q\varphi$ 可知污泥絮体在 M 点的电势能比在 N 点的电势能大，污泥絮体从 M 点移到 N 点，电势能减小，电场力对其做正功，故 AC 正确；

B．根据电场线的疏密程度可知 N 点的电场强度比 P 点的小，故 B 错误；

D．M 点和 P 点在同一等势面上，则污泥絮体在 M 点的电势能与在 P 点的电势能相等，结合 AC 选项分析可知污泥絮体在 P 点的电势能比其在 N 点的大，故 D 错误。

故选 AC。
}

\item
%% number: 9
%% paperName: 2024年广东高考第 9 题 6 分
%% typeId: 2
%% type: 多选题
%% chapter: 曲线运动
%% point: 万有引力定律
%% method: 无
%% score: 6
%% degree: 550
%% duplicateId: 0
%% body:
如图所示，探测器及其保护背罩通过弹性轻绳连接降落伞。在接近某行星表面时以 $60 \Ums$ 的速度竖直匀速下落。此时启动“背罩分离”，探测器与背罩断开连接，背罩与降落伞保持连接。已知探测器质量为 $1000 \Ukg$，背罩质量为 $50 \Ukg$，该行星的质量和半径分别为地球的 $\frac{1}{10}$ 和 $\frac{1}{2}$。地球表面重力加速度大小取 $g = 10 \Umsq$。忽略大气对探测器和背罩的阻力。下列说法正确的有\xzanswer{AC}

\onepicture[w=4.50cm, align=right]{figs/09.png}

\fourchoices[answer=AC]
{该行星表面的重力加速度大小为 $4 \Umsq$}
{该行星的第一宇宙速度为 $7.9 \Ukms$}
{“背罩分离”后瞬间，背罩的加速度大小为 $80 \Umsq$}
{“背罩分离”后瞬间，探测器所受重力对其做功的功率为 $30 \UkW$}

%% answer: AC
%% memo:
\memoanswer{
A．在星球表面，根据 $G\frac{Mm}{R^{2}}=mg$

可得 $g=\frac{GM}{R^{2}}$

行星的质量和半径分别为地球的 $\frac{1}{10}$ 和 $\frac{1}{2}$。地球表面重力加速度大小取 $g = 10 \Umsq$，可得该行星表面的重力加速度大小 $g' = 4 \Umsq$

故 A 正确；

B．在星球表面上空，根据万有引力提供向心力

$G\frac{Mm}{R^{2}}=m\frac{v^{2}}{R}$

可得星球的第一宇宙速度

$v=\sqrt{\frac{GM}{R}}$

行星的质量和半径分别为地球的 $\frac{1}{10}$ 和 $\frac{1}{2}$，可得该行星的第一宇宙速度

$v_{\text{行}}=\frac{\sqrt{5}}{5}v_{\text{地}}$

地球的第一宇宙速度为 $7.9 \Ukms$，所以该行星的第一宇宙速度

$v_{\text{行}}=\frac{\sqrt{5}}{5}\times7.9 \Ukms$

故 B 错误；

C．“背罩分离”前，探测器及其保护背罩和降落伞整体做匀速直线运动，对探测器受力分析，可知探测器与保护背罩之间的作用力 $F = mg' = 4000 \UN$

“背罩分离”后，背罩所受的合力大小为 $4000 \UN$，对背罩，根据牛顿第二定律 $F = m'a$

解得 $a = 80 \Umsq$

故 C 正确；

D．“背罩分离”后瞬间探测器所受重力对其做功的功率

$P = mg'v = 1000 \times 4 \times 60 \UW = 240 \UkW$

故 D 错误。

故选 AC。
}

\item
%% number: 10
%% paperName: 2024年广东高考第 10 题 6 分
%% typeId: 2
%% type: 多选题
%% chapter: 运动和力
%% point: 牛顿定律的应用
%% method: 无
%% score: 6
%% degree: 450
%% duplicateId: 0
%% body:
如图所示，光滑斜坡上，可视为质点的甲、乙两个相同滑块，分别从 $H_{\text{甲}}$、$H_{\text{乙}}$ 高度同时由静止开始下滑。斜坡与水平面在 O 处平滑相接，滑块与水平面间的动摩擦因数为 $\mu$，乙在水平面上追上甲时发生弹性碰撞。忽略空气阻力。下列说法正确的有\xzanswer{ABD}

\onepicture[w=6.00cm, align=right]{figs/10.png}

\fourchoices[answer=ABD]
{甲在斜坡上运动时与乙相对静止}
{碰撞后瞬间甲的速度等于碰撞前瞬间乙的速度}
{乙的运动时间与 $H_{\text{乙}}$ 无关}
{甲最终停止位置与 O 处相距 $\frac{H_{\text{乙}}}{\mu}$}

%% answer: ABD
%% memo:
\memoanswer{
A．两滑块在光滑斜坡上加速度相同，同时由静止开始下滑，则相对速度为 0，故 A 正确；

B．两滑块滑到水平面后均做匀减速运动，由于两滑块质量相同，且发生弹性碰撞，可知碰后两滑块交换速度，即碰撞后瞬间甲的速度等于碰撞前瞬间乙的速度，故 B 正确；

C．设斜面倾角为 $\theta$，乙下滑过程有

$H_{\text{乙}} = \frac{1}{2}g\sin\theta t_{1}^{2}$

在水平面运动一段时间 $t_{2}$ 后与甲相碰，碰后以甲碰前速度做匀减速运动 $t_{3}$，乙运动的时间为

$t = t_{1} + t_{2} + t_{3}$

由于 $t_{1}$ 与 $H_{\text{乙}}$ 有关，则总时间与 $H_{\text{乙}}$ 有关，故 C 错误；

D．乙下滑过程有

$mgH_{\text{乙}} = \frac{1}{2}mv_{\text{乙}}^{2}$

由于甲和乙发生弹性碰撞，交换速度，则可知甲最终停止位置与不发生碰撞时乙最终停止的位置相同；则如果不发生碰撞，乙在水平面运动到停止有 $v_{\text{乙}}^{2}=2\mu gx$

联立可得 $x=\frac{H_{\text{乙}}}{\mu}$

即发生碰撞后甲最终停止位置与 O 处相距 $\frac{H_{\text{乙}}}{\mu}$。故 D 正确。

故选 ABD。
}

\item
%% number: 11
%% paperName: 2024年广东高考第 11 题 7 分
%% typeId: 4
%% type: 实验题
%% chapter: 运动和力
%% point: 验证牛顿第二定律
%% method: 无
%% score: 7
%% degree: 750
%% duplicateId: 0
%% body:
下列是《普通高中物理课程标准》中列出的三个必做实验的部分步骤，请完成实验操作和计算。

（1）图甲是“探究加速度与物体受力、物体质量的关系”实验装置示意图。图中木板右端垫高的目的是\tkanswer[2.5cm]{平衡摩擦力}。图乙是实验得到纸带的一部分，每相邻两计数点间有四个点未画出。相邻计数点的间距已在图中给出。打点计时器电源频率为 $50 \UHz$，则小车的加速度大小为\tkanswer[1.5cm]{2.86}$\Umsq$（结果保留 3 位有效数字）。

（2）在“长度的测量及其测量工具的选用”实验中，某同学用 50 分度的游标卡尺测量一圆柱体的长度，示数如图丙所示，图丁为局部放大图，读数为\tkanswer[1.5cm]{4.108}$\Ucm$。

（3）在“用双缝干涉实验测量光的波长”实验调节过程中，在光具座上安装光源、遮光筒和光屏。遮光筒不可调节。打开并调节\tkanswer[1.5cm]{光源}。使光束沿遮光筒的轴线把光屏照亮。取下光屏，装上单缝、双缝和测量头。调节测量头，并缓慢调节单缝的角度直到目镜中观察到\tkanswer[2cm]{干涉条纹}。

\onepicture[w=9.00cm, align=right]{figs/11a.png}

\onepicture[w=9.00cm, align=right]{figs/11b.png}

%% answer:
\jdanswer{
\begin{enumerate}
	\item
	平衡摩擦力，$2.86$

	\item
	$4.108$

	\item
	光源，干涉条纹

\end{enumerate}
}
%% memo:
\memoanswer{
（1）木板右端抬高的目的是平衡摩擦力；

小车的加速度大小

$a=\frac{(16.29+13.43+10.59-7.72-4.88-2.01)\times10^{-2}}{9\times0.1^{2}} \Umsq=2.86 \Umsq$

（2）游标卡尺读数为 $4.1 \Ucm + 4 \times 0.02 \Umm = 4.108 \Ucm$

（3）在“用双缝干涉实验测量光的波长”实验调节过程中，在光具座上安装光源、遮光筒和光屏，遮光筒不可调节，打开并调节光源，使光束沿遮光筒的轴线把光屏照亮，取下光屏，装上单缝、双缝和测量头，调节测量头，并缓慢调节单缝的角度直到目镜中观察到干涉条纹。
}

\item
%% number: 12
%% paperName: 2024年广东高考第 12 题 9 分
%% typeId: 4
%% type: 实验题
%% chapter: 电路
%% point: 电路实验
%% method: 无
%% score: 9
%% degree: 450
%% duplicateId: 0
%% body:
某科技小组模仿太阳能发电中的太阳光自动跟踪系统，制作光源跟踪演示装置，实现太阳能电池板方向的调整，使电池板正对光源。图甲是光照方向检测电路。所用器材有：电源 $E$（电动势 $3 \UV$）；电压表（V$_{1}$）和（V$_{2}$）（量程均有 $3 \UV$ 和 $15 \UV$，内阻均可视为无穷大）；滑动变阻器 $R$；两个相同的光敏电阻 $R_{G1}$ 和 $R_{G2}$；开关 S；手电筒；导线若干。图乙是实物图。图中电池板上垂直安装有半透明隔板，隔板两侧装有光敏电阻，电池板固定在电动机转轴上。控制单元与检测电路的连接未画出。控制单元对光照方向检测电路无影响。请完成下列实验操作和判断。

（1）电路连接。

图乙中已正确连接了部分电路，请完成虚线框中滑动变阻器 $R$、电源 $E$、开关 S 和电压表 V$_{1}$ 间的实物图连线\tkanswer[2cm]{}。

（2）光敏电阻阻值与光照强度关系测试。

①将图甲中 $R$ 的滑片置于\tkanswer[1cm]{b}端。用手电筒的光斜照射到 $R_{G1}$ 和 $R_{G2}$，使 $R_{G1}$ 表面的光照强度比 $R_{G2}$ 表面的小。

②闭合 S，将 $R$ 的滑片缓慢滑到某一位置。V$_{1}$ 的示数如图丙所示，读数 $U_{1}$ 为\tkanswer[1cm]{1.60}$\UV$，V$_{2}$ 的示数 $U_{2}$ 为 $1.17 \UV$。由此可知，表面光照强度较小的光敏电阻的阻值\tkanswer[1cm]{较大}（填“较大”或“较小”）。

③断开 S。

（3）光源跟踪测试。

①将手电筒的光从电池板上方斜照射到 $R_{G1}$ 和 $R_{G2}$。

②闭合 S，并启动控制单元。控制单元检测并比较两光敏电阻的电压，控制电动机转动。此时两电压表的示数 $U_{1} < U_{2}$，图乙中的电动机带动电池板\tkanswer[1.5cm]{逆时针}（填“逆时针”或“顺时针”）转动，直至\tkanswer[1.5cm]{$U_{1} = U_{2}$}时停止转动，电池板正对手电筒发出的光。

\onepicture[w=4.20cm, align=right]{figs/12a.png}

\onepicture[w=7.00cm, align=right]{figs/12b.png}

\onepicture[w=3.80cm, align=right]{figs/12c.png}

%% answer:
\jdanswer{
\begin{enumerate}
	\item
	电路连线见解析图。

	\item
	$b$，$1.60$，较大。

	\item
	逆时针，$U_{1} = U_{2}$。

\end{enumerate}
}
%% memo:
\memoanswer{
（1）电路连线如图。

\onepicture[w=6.50cm]{figs/12_an.jpg}

（2）①为了保证电路的安全，实验开始前要将 $R$ 的滑片置于阻值最大处，即置于 b 端。

②电压表 V$_{1}$ 量程为 $3 \UV$，最小刻度为 $0.1 \UV$，则读数为 $1.60 \UV$；电压表 V$_{1}$ 比电压表 V$_{2}$ 的示数大，说明 $R_{G1} > R_{G2}$，由此可知表面光照强度较小的光敏电阻的阻值较大。

（3）电压表的示数 $U_{1} < U_{2}$，说明 $R_{G1}$ 表面的光照强度比 $R_{G2}$ 表面的大，因此电动机带动电池板逆时针转动，直至 $U_{1} = U_{2}$ 时停止转动，电池板正对手电筒发出的光。
}

\item
%% number: 13
%% paperName: 2024年广东高考第 13 题 9 分
%% typeId: 6
%% type: 计算题
%% chapter: 气体性质
%% point: 查理定律
%% method: 无
%% score: 9
%% degree: 550
%% duplicateId: 0
%% body:
差压阀可控制气体进行单向流动，广泛应用于减震系统。如图所示，A、B 两个导热良好的气缸通过差压阀连接，A 内轻质活塞的上方与大气连通，B 内气体体积不变。当 A 内气体压强减去 B 内气体压强大于 $\Delta p$ 时差压阀打开，A 内气体缓慢进入 B 中；当该差值小于或等于 $\Delta p$ 时差压阀关闭。当环境温度 $T_{1} = 300 \UK$ 时，A 内气体体积 $V_{A1} = 4.0 \times 10^{2} \Umc$，B 内气体压强 $p_{B1}$ 等于大气压强 $p_{0}$，已知活塞的横截面积 $S = 0.10 \Umq$，$\Delta p = 0.11p_{0}$，$p_{0} = 1.0 \times 10^{5} \UPa$，重力加速度大小取 $g = 10 \Umsq$，A、B 内的气体可视为理想气体，忽略活塞与气缸间的摩擦、差压阀与连接管内的气体体积不计。当环境温度降到 $T_{2} = 270 \UK$ 时：
\begin{enumerate}
	\item
	求 B 内气体压强 $p_{B2}$；
	\item
	求 A 内气体体积 $V_{A2}$；
	\item
	在活塞上缓慢倒入铁砂，若 B 内气体压强回到 $p_{0}$ 并保持不变，求已倒入铁砂的质量 $m$。
\end{enumerate}
\onepicture[w=5.50cm, align=right]{figs/13.png}
%% answer:
\jdanswer{
\begin{enumerate}
	\item
	$9 \times 10^{4} \UPa$

	\item
	$3.6 \times 10^{2} \Umc$

	\item
	$1.1 \times 10^{2} \Ukg$

\end{enumerate}
}
%% memo:
\memoanswer{
（1）（2）假设温度降低到 $T_{2}$ 时，差压阀没有打开，A、B 两个气缸导热良好，B 内气体做等容变化，初态 $p_{B1}=p_{0}$，$T_{1}=300 \UK$

末态 $T_{2}=270 \UK$

根据 $\frac{p_{B1}}{T_{1}}=\frac{p_{B2}}{T_{2}}$

代入数据可得 $p_{B2}=9 \times 10^{4} \UPa$

A 内气体做等压变化，压强保持不变，初态 $V_{A1}=4.0 \times 10^{2} \Umc$，$T_{1}=300 \UK$

末态 $T_{2}=270 \UK$

根据 $\frac{V_{A1}}{T_{1}}=\frac{V_{A2}}{T_{2}}$

代入数据可得 $V_{A2}=3.6 \times 10^{2} \Umc$

由于 $p_{0}-p_{B2}<\Delta p$

假设成立，即 $p_{B2}=9 \times 10^{4} \UPa$

（3）恰好稳定时，A 内气体压强为 $p_{A}'=p_{0}+\frac{mg}{S}$

B 内气体压强 $p_{B}'=p_{0}$

此时差压阀恰好关闭，所以有 $p_{A}'-p_{B}'=\Delta p$

代入数据联立解得 $m = 1.1 \times 10^{2} \Ukg$。
}

\item
%% number: 14
%% paperName: 2024年广东高考第 14 题 13 分
%% typeId: 6
%% type: 计算题
%% chapter: 运动和力
%% point: 牛顿第一定律
%% method: 无
%% score: 13
%% degree: 550
%% duplicateId: 0
%% body:
汽车的安全带和安全气囊是有效保护乘客的装置。
\begin{enumerate}
	\item
	安全带能通过感应车的加速度自动锁定，其原理的简化模型如图甲所示。在水平路面上刹车的过程中，敏感球由于惯性沿底座斜面上滑直到与车达到共同的加速度 $a$，同时顶起敏感臂，使之处于水平状态，并卡住卷轴外齿轮，锁定安全带。此时敏感臂对敏感球的压力大小为 $F_{N}$，敏感球的质量为 $m$，重力加速度为 $g$。忽略敏感球受到的摩擦力。求斜面倾角的正切值 $\tan\theta$。
	\item
	如图乙所示，在安全气囊的性能测试中，可视为质点的头锤从离气囊表面高度为 $H$ 处做自由落体运动，与正下方的气囊发生碰撞。以头锤到气囊表面为计时起点，气囊对头锤竖直方向作用力 $F$ 随时间 $t$ 的变化规律，可近似用图丙所示的图像描述。已知头锤质量 $M = 30 \Ukg$，$H = 3.2 \Um$，重力加速度大小取 $g = 10 \Umsq$。求：①碰撞过程中 $F$ 的冲量大小和方向；②碰撞结束后头锤上升的最大高度。
\end{enumerate}
\onepicture[w=4.00cm, align=right]{figs/14a.png}

\onepicture[w=8.00cm, align=right]{figs/14b.png}
%% answer:
\jdanswer{
\begin{enumerate}
	\item
	$\tan\theta = \frac{ma}{mg + F_{N}}$

	\item
	① $330 \UNs$，方向竖直向上；② $0.2 \Um$

\end{enumerate}
}
%% memo:
\memoanswer{
（1）敏感球受向下的重力 $mg$ 和敏感臂向下的压力 $F_{N}$ 以及斜面的支持力 $N$，则由牛顿第二定律可知

$\left(mg+F_{N}\right)\tan\theta=ma$

解得 $\tan\theta=\frac{ma}{mg+F_{N}}$

（2）①由图像可知碰撞过程中 $F$ 的冲量大小

$I_{F}=\frac{1}{2}\times0.1\times6600 \UNs=330 \UNs$

方向竖直向上；

②头锤落到气囊上时的速度

$v_{0}=\sqrt{2gH}=8 \Ums$

与气囊作用过程由动量定理（向上为正方向）

$I_{F}-mgt=mv-(-mv_{0})$

解得 $v=2 \Ums$

则上升的最大高度

$h=\frac{v^{2}}{2g}=0.2 \Um$。
}

\item
%% number: 15
%% paperName: 2024年广东高考第 15 题 16 分
%% typeId: 6
%% type: 计算题
%% chapter: 磁场
%% point: 带电粒子在组合场中的运动
%% method: 无
%% score: 16
%% degree: 350
%% duplicateId: 0
%% body:
如图甲所示，两块平行正对的金属板水平放置，板间加上如图乙所示幅值为 $U_{0}$、周期为 $t_{0}$ 的交变电压。金属板左侧存在一水平向右的恒定匀强电场，右侧分布着垂直纸面向外的匀强磁场。磁感应强度大小为 $B$。一带电粒子在 $t = 0$ 时刻从左侧电场某处由静止释放，在 $t = t_{0}$ 时刻从下板左端边缘位置水平向右进入金属板间的电场内，在 $t = 2t_{0}$ 时刻第一次离开金属板间的电场、水平向右进入磁场，并在 $t = 3t_{0}$ 时刻从下板右端边缘位置再次水平进入金属板间的电场。已知金属板的板长是板间距离的 $\frac{\pi}{3}$ 倍，粒子质量为 $m$。忽略粒子所受的重力和场的边缘效应。
\begin{enumerate}
	\item
	判断带电粒子的电性并求其所带的电荷量 $q$；
	\item
	求金属板的板间距离 $D$ 和带电粒子在 $t = t_{0}$ 时刻的速度大小 $v$；
	\item
	求从 $t = 0$ 时刻开始到带电粒子最终碰到上金属板的过程中，电场力对粒子做的功 $W$。
\end{enumerate}
\onepicture[w=8.00cm, align=right]{figs/15.png}
%% answer:
\jdanswer{
\begin{enumerate}
	\item
	正电；$q = \frac{\pi m}{Bt_{0}}$

	\item
	$D = \sqrt{\frac{3\pi t_{0}U_{0}}{8B}}$，$v = \pi\sqrt{\frac{\pi U_{0}}{24Bt_{0}}}$

	\item
	$W = \frac{\pi mU_{0}(\pi^{2} + 16)}{48Bt_{0}}$

\end{enumerate}
}
%% memo:
\memoanswer{
（1）根据带电粒子在右侧磁场中的运动轨迹结合左手定则可知，粒子带正电；粒子在磁场中运动的周期为 $T = 2t_{0}$

根据 $T=\frac{2\pi m}{qB}$

则粒子所带的电荷量 $q=\frac{\pi m}{Bt_{0}}$

（2）若金属板的板间距离为 $D$，则板长 $\frac{\pi D}{3}$，粒子在板间运动时

$\frac{\pi D}{3}=v t_{0}$

出电场时竖直速度为零，则竖直方向

$y=2\times\frac{1}{2}\frac{U_{0}q}{Dm}(0.5t_{0})^{2}$

在磁场中时

$qvB=m\frac{v^{2}}{r}$

其中

$y=2r=\frac{2mv}{qB}$

联立解得

$D=\sqrt{\frac{3\pi t_{0}U_{0}}{8B}}$，$v=\pi\sqrt{\frac{\pi U_{0}}{24Bt_{0}}}$

（3）带电粒子在电场和磁场中的运动轨迹如图，由（2）的计算可知金属板的板间距离 $D = 3r$。

\onepicture[w=6.50cm]{figs/15_an.jpg}

则粒子在 $3t_{0}$ 时刻再次进入中间的偏转电场，在 $4t_{0}$ 时刻进入左侧的电场做减速运动速度为零后反向加速，在 $6t_{0}$ 时刻再次进入中间的偏转电场，$6.5t_{0}$ 时刻碰到上极板，因粒子在偏转电场中运动时，在时间 $t_{0}$ 内电场力做功为零，在左侧电场中运动时，往返一次电场力做功也为零，可知整个过程中只有开始进入左侧电场时电场力做功和最后 $0.5t_{0}$ 时间内电场力做功，则

$W=\frac{1}{2}mv^{2}+Eq\times\frac{D}{3}=\frac{\pi^{3}mU_{0}}{48Bt_{0}}+\frac{\pi mU_{0}}{3Bt_{0}}=\frac{\pi mU_{0}(\pi^{2}+16)}{48Bt_{0}}$。
}

\end{enumerate}

\end{document}
